\section{问题分析}

\subsection{问题一分析}
问题一的核心是建立可靠的数据基础。先检查取值范围、重复记录和缺失比例，再根据指标含义选择中位数填补、截尾或保留缺失标记等方法，最后分别处理效益型和成本型指标。

\subsection{问题二分析}
问题二需要把多维指标映射为可解释的评价结果。可结合题意、专家判断或客观赋权确定权重，再通过线性加权、TOPSIS 等方法计算综合得分，并把资源配置写成带预算和容量约束的优化问题。

\subsection{问题三分析}
问题三关注结论是否依赖某一组参数。可对权重、预算或观测误差进行扰动，比较排序、目标函数与关键决策变量的变化，并据此划定模型的稳定适用区间。
